% ID: SEL-D0046-OA-E02
% ESTADO: cotejada. Segunda pasada visual y borrador compilado revisados.
% Fuente: Selectividad/sel_2014_matematicas_aplicadas/Matematicas_CCSS_JUN.pdf
% SHA-256: df7375f8183eafb6f45e9c65083e68d528ffabe07d66717b0ccf0f020c5b064f
% Página: 1; opción A; ejercicio original 2.

La función de beneficios $f$, en miles de euros, de una empresa depende de la cantidad invertida $x$, en miles de euros, en un determinado proyecto de innovación y viene dada por $f(x)=-2x^2+36x+138$, $x\geq0$.
\begin{enumerate}[label=\alph*)]
\item \textbf{(1 punto)} Determine la inversión que maximiza el beneficio de la empresa y calcule dicho beneficio óptimo.
\item \textbf{(0.5 puntos)} Calcule $f'(7)$ e interprete el signo del resultado.
\item \textbf{(1 punto)} Dibuje la función de beneficios $f(x)$. \textquestiondown{}Para qué valor o valores de la inversión, $x$, el beneficio es de 138 mil euros?
\end{enumerate}
